El servicio municipal de saneamiento de una gran ciudad debe inspeccionar la red principal de colectores. Cada tramo debe revisarse una vez, y el equipo de inspección quiere planificar un recorrido que empiece y termine en el mismo punto minimizando la distancia total recorrida.

La red se modeliza mediante un grafo no dirigido y ponderado, donde los vértices representan puntos de acceso a los colectores y las aristas representan tramos, con sus longitudes en metros.

Los datos vienen dados por la siguiente matriz de adyacencia (0 indica ausencia de conexión directa):

\[
\begin{array}{c|cccccccccc}
 & A & B & C & D & E & F & G & H & I & J \\
\hline
A & 0 & 90 & 120 & 110 & 0 & 0 & 0 & 0 & 0 & 0 \\
B & 90 & 0 & 80 & 0 & 100 & 0 & 0 & 0 & 0 & 0 \\
C & 120 & 80 & 0 & 70 & 0 & 95 & 0 & 0 & 0 & 0 \\
D & 110 & 0 & 70 & 0 & 0 & 0 & 85 & 0 & 0 & 0 \\
E & 0 & 100 & 0 & 0 & 0 & 75 & 0 & 60 & 0 & 0 \\
F & 0 & 0 & 95 & 0 & 75 & 0 & 60 & 0 & 50 & 0 \\
G & 0 & 0 & 0 & 85 & 0 & 60 & 0 & 0 & 0 & 80 \\
H & 0 & 0 & 0 & 0 & 60 & 0 & 0 & 0 & 70 & 0 \\
I & 0 & 0 & 0 & 0 & 0 & 50 & 0 & 70 & 0 & 65 \\
J & 0 & 0 & 0 & 0 & 0 & 0 & 80 & 0 & 65 & 0 \\
\end{array}
\]

Calcula la distancia total mínima y propón un recorrido óptimo. (Justifica todos los pasos y explica claramente el procedimiento seguido).
